% Copyright (c) 2026 Twin Vector Labs LLC.
% All rights reserved.
\documentclass[11pt]{article}

\usepackage[margin=1in]{geometry}
\usepackage{amsmath,amssymb,amsthm}
\usepackage{mathtools}
\usepackage{enumitem}
\usepackage[colorlinks=true,linkcolor=blue,urlcolor=blue]{hyperref}

\newtheorem{theorem}{Theorem}
\newtheorem{proposition}{Proposition}
\newtheorem{corollary}{Corollary}
\theoremstyle{definition}
\newtheorem{remark}{Remark}

\newcommand{\C}{\mathbb{C}}
\newcommand{\Z}{\mathbb{Z}}
\newcommand{\rep}[1]{\mathbf{#1}}      % a representation, e.g. \rep{3}
\newcommand{\triv}{\mathbf{1}}         % trivial rep / singlet
\newcommand{\sgn}{\mathrm{sgn}}
\newcommand{\Sym}{\operatorname{Sym}}
\newcommand{\sigmacol}{\sigma}         % color charge functional

\title{\bf Why the Proton Is Irreducibly Tripartite:\\
the color-singlet obstruction to bipartite qudit merges}
\author{Tessera --- cobordism subsystem (\#382 / \#396)}
\date{}

\begin{document}
\maketitle

\begin{abstract}
We attempted to build a proton by \emph{composing bipartite merges}: relax a
cobordism whose boundary is two color states $\psi_A,\psi_B$, read the emergent
result $\psi_{AB}$, then feed it as a boundary input of a second merge with a
third state $\psi_C$. The emergent result never reaches the color singlet --- it
stays color-charged ($|\sigmacol|\approx 0.67$) at every step. This note shows
that this is \emph{not} a numerical or topological defect of the construction but
a representation-theoretic theorem: the color singlet of $SU(3)$ is the
three-index alternating invariant $\varepsilon_{ijk}$, which is \emph{absent}
from the two-fold tensor product $\rep{3}\otimes\rep{3}$ and first appears in
$\rep{3}\otimes\rep{3}\otimes\rep{3}$. A baryon singlet built from fundamentals
is therefore irreducibly an $n=3$ object; no single bipartite fusion of two
color triplets can produce it. The qudit generalization makes the statement
sharp: for $SU(d)$ the baryon arity equals $d$, so a \emph{qubit} ($d=2$) baryon
\emph{is} bipartite (the spin singlet $\varepsilon_{ij}$), while a \emph{qutrit}
color baryon ($d=3$) is not. The simplicial resolution is the trivalent junction
$W_{ABC}$ that carries $\varepsilon_{ijk}$ directly.
\end{abstract}

\section{The question}

The merge primitive is a Lorentzian Regge cobordism $W$ with prescribed boundary
$\partial W$: we pin the boundary color states, relax the interior edge lengths
to a stationary point of the dual Regge action subject to a realizability
constraint, and read the emergent result off the relaxed geometry. Composition
is by \emph{cobordism}, not welding: the emergent result-state of one merge is
the boundary input of the next.

Concretely we merged three color-neutral $q\bar q$ pairs
$A=(1,-1,0)$, $B=(1,0,-1)$, $C=(0,1,-1)$ pairwise:
\[
  \text{merge}(A,B)\to AB,\qquad \text{merge}(AB,C)\to ABC .
\]
With the exact period read-out (\texttt{residualForPeriods} for the pinned
inputs, \texttt{cyclePeriods} for the emergent result block) we measure the
emergent color charge $\sigmacol(\psi)=\sum_i\psi_i$:
\[
  |\sigmacol_{AB}|\approx 0.70,\qquad |\sigmacol_{ABC}|\approx 0.67,
  \qquad\text{target singlet } |\sigmacol|=0 .
\]
The result is color-charged at every step. \emph{Is this a bug or a theorem?}
It is a theorem.

\section{Color, invariant tensors, and singlets}

A quark transforms in the fundamental $\rep{3}$ of $SU(3)_{\text{color}}$, an
antiquark in $\bar{\rep{3}}$. A physical (observable) hadron must be
\emph{colorless}: a singlet, i.e.\ an invariant under the color group. Building
hadrons is therefore the problem of constructing $SU(3)$-invariant tensors out of
copies of $\rep{3}$ and $\bar{\rep 3}$.

\begin{theorem}[First fundamental theorem of invariant theory for $SL(3)$]
Every $SU(3)$-invariant tensor in copies of $\rep{3}$ and $\bar{\rep 3}$ is a
polynomial in the two elementary invariants
\[
  \delta^{i}{}_{j}\quad(\text{the contraction }\rep{3}\otimes\bar{\rep 3}\to\triv),
  \qquad
  \varepsilon_{ijk},\ \varepsilon^{ijk}\quad(\text{the alternating }\Lambda^3).
\]
\end{theorem}

The two invariants are exactly the two ways to make a colorless object:
\begin{itemize}[nosep]
  \item $\delta^{i}{}_{j}$ pairs a quark with an antiquark: the \textbf{meson}
  $\bar q_i\,q^i$. It is \emph{bilinear} --- a colorless meson is a two-body
  (bipartite) object.
  \item $\varepsilon_{ijk}$ contracts \emph{three} quarks:
  the \textbf{baryon} $\varepsilon_{ijk}\,q^i q^j q^k$ (the proton). It is
  \emph{trilinear} and totally antisymmetric.
\end{itemize}

The decisive fact is what is \emph{missing}: there is no invariant bilinear in two
fundamentals.

\begin{proposition}[No singlet at two bodies]
\label{prop:nosinglet}
\[
  \rep{3}\otimes\rep{3} \;=\; \rep{6}\,\oplus\,\bar{\rep 3},
  \qquad\text{and neither summand is the trivial representation }\triv .
\]
Hence the two-quark Hilbert space $\rep 3\otimes\rep 3$ contains \emph{no} color
singlet. Equivalently, the multiplicity of $\triv$ in $\rep 3^{\otimes 2}$ is $0$.
\end{proposition}

The singlet first appears one body later:
\begin{proposition}[The singlet is born at three bodies]
\label{prop:three}
\[
  \rep 3^{\otimes 2}=\rep 6\oplus\bar{\rep 3}
  \quad(\dim 9 = 6+3),
  \qquad
  \rep 3^{\otimes 3}=\rep{10}\oplus\rep 8\oplus\rep 8\oplus\triv
  \quad(\dim 27 = 10+8+8+1).
\]
The trivial representation occurs with multiplicity $0,0,1$ in
$\rep 3^{\otimes 1},\rep 3^{\otimes 2},\rep 3^{\otimes 3}$; the unique copy in
$\rep 3^{\otimes 3}$ is spanned by $\varepsilon_{ijk}$.
\end{proposition}

Proposition~\ref{prop:nosinglet} is the heart of the matter: \emph{the proton
state is not an element of any two-body color Hilbert space}. A bipartite fusion
$V_A\otimes V_B\to V_R$ of two color triplets can only land in
$\rep 6\oplus\bar{\rep 3}$; its output spectrum simply does not contain the
singlet to which we are trying to relax. Compare the colorless meson: it lives in
$\rep 3\otimes\bar{\rep 3}=\triv\oplus\rep 8$, which \emph{does} contain
$\triv$ --- which is exactly why a $q\bar q$ pair (our carriable input) is a
legitimate two-body colorless object, while a proton is not.

\section{The associativity subtlety: why a \emph{sequence} of merges still fails}

One might object that the singlet \emph{is} reachable pairwise through
associativity:
\[
  (\rep 3\otimes\rep 3)\otimes\rep 3
  =(\rep 6\oplus\bar{\rep 3})\otimes\rep 3
  =\underbrace{(\rep 6\otimes\rep 3)}_{\rep{10}\oplus\rep 8}
   \;\oplus\;
   \underbrace{(\bar{\rep 3}\otimes\rep 3)}_{\triv\oplus\rep 8}.
\]
The singlet appears \emph{only} in the $\bar{\rep 3}\otimes\rep 3$ channel: one
must first project the diquark onto its \emph{antisymmetric} part
$\bar{\rep 3}=\varepsilon_{ijk}q^j q^k$ and discard the symmetric sextet
$\rep 6$. A bipartite merge that \emph{projects onto the antisymmetric channel}
could then complete the singlet with the third quark.

But the emergent merge does no such projection. Relaxing the cobordism produces a
stationary geometry carrying the \emph{generic} product $\rep 6\oplus\bar{\rep 3}$;
the symmetric $\rep 6$ component is not annihilated, and it contaminates the
second merge --- the result stays colored ($|\sigmacol_{AB}|\approx0.70$). To
kill the $\rep 6$ one must impose the \emph{simultaneous} antisymmetry of all
three color indices, $\varepsilon_{ijk}$, at a \emph{single} vertex. That object
is a trivalent junction carrying the alternating (sign) representation --- not any
composition of generic pairwise merges.

\begin{remark}
This is the precise sense in which the proton is ``irreducibly tripartite'': the
defining invariant $\varepsilon_{ijk}\in\Lambda^3(\C^3)^*$ is a totally
antisymmetric \emph{three}-form. It does not factor as a composite of bilinear
invariants except by first manufacturing the antisymmetric diquark $\bar{\rep 3}$
--- and that projection is itself the statement that the three indices must meet
antisymmetrically. There is no associative bracketing that turns a \emph{generic}
two-body fusion into a singlet.
\end{remark}

\section{The qudit generalization (the sharp answer)}

The argument is not special to $SU(3)$; it pins the baryon \emph{arity} exactly.
For $SU(d)$ acting on a qudit (the fundamental $\C^d$), the invariant tensors are
again generated by $\delta^{i}{}_{j}$ and the Levi-Civita symbol
$\varepsilon_{i_1\cdots i_d}$, now with \emph{$d$ indices}.

\begin{corollary}[Baryon arity $=d$]
A singlet built from fundamentals of $SU(d)$ requires exactly $d$ of them:
$\varepsilon_{i_1\cdots i_d}\,q^{i_1}\!\cdots q^{i_d}$. The multiplicity of
$\triv$ in $\rep d^{\otimes n}$ is $0$ for $0<n<d$ and first becomes nonzero at
$n=d$.
\end{corollary}

This is the direct answer to ``why can't bipartite qubit/qudit interactions reach
the proton'':
\begin{itemize}[nosep]
  \item \textbf{Qubit, $d=2$.} $\varepsilon_{ij}$ has \emph{two} indices:
  $\rep 2\otimes\rep 2=\rep 3\oplus\triv$ already contains a singlet (the spin-$0$
  state $\tfrac{1}{\sqrt2}(\!\uparrow\downarrow-\downarrow\uparrow)$). A
  ``baryon'' of qubits \emph{is} bipartite. Bipartite merges suffice here ---
  which is exactly why qubit intuition misleads.
  \item \textbf{Qutrit / color, $d=3$.} $\varepsilon_{ijk}$ has \emph{three}
  indices: $\rep 3\otimes\rep 3=\rep 6\oplus\bar{\rep 3}$ has no singlet. The
  proton is a qutrit baryon and is irreducibly \emph{tripartite}.
\end{itemize}
The proton's three-quark structure is, in one line, the statement $d=3$: the
color singlet is the $d$-fold alternating invariant, and $d=3$ for color.

\section{The simplicial model and the measurement}

In the cobordism subsystem the color register is the homology of the holed
register surface: $b_1=2$ on the $\sigmacol=0$ hyperplane, carrying the
\emph{standard} ($2$-dimensional) representation of $S_3=\mathrm{Weyl}(SU(3))$
permuting the three colors, with
\[
  \C^3 \;=\; \underbrace{\triv}_{\sigmacol\neq0\ \text{direction}}
            \;\oplus\;
            \underbrace{V_{\text{std}}}_{\sigmacol=0\ \text{register}} .
\]
\textbf{Confinement} is the topological statement that the $\sigmacol\neq0$
(trivial-rep) direction is \emph{not carried} by the register: a net color charge
has no realizable harmonic representative, while a $\sigmacol=0$ configuration
does. The proton singlet is represented by the definite-triality state
\[
  \psi_{\text{proton}}=(1,\omega,\omega^2),\qquad \omega=e^{2\pi i/3},
  \qquad 1+\omega+\omega^2=0 ,
\]
the $\Z_3\subset S_3$ Fourier mode: it is the discrete shadow of the alternating
invariant $\varepsilon_{ijk}$ (an eigenvector of the cyclic color rotation, i.e.\
a state of definite triality, lying in the $\sigmacol=0$ register).

A bipartite merge is a cobordism $W$ with $\partial W=\psi_A\sqcup\psi_B$ and an
emergent result block $R$; we read $\sigmacol_R=\sum_i (\psi_R)_i$ from the
relaxed metric via \texttt{cyclePeriods}. Consistent with
Proposition~\ref{prop:nosinglet}, the emergent two-input result does not land in
the singlet direction:
\[
  |\sigmacol_{AB}|\approx 0.70,\qquad
  |\sigmacol_{ABC}|\approx 0.67 \quad(\text{sequence}),
\]
against the singlet target $|\sigmacol|=0$. Two control facts confirm the reading
is faithful and the obstruction structural:
\begin{itemize}[nosep]
  \item The result is read with the \emph{exact} period machinery
  (\texttt{residualForPeriods}/\texttt{cyclePeriods}), not the soft edge-loop
  encoding; the $\approx10^{-2}$ neutral realizability floor is intrinsic to the
  shared-register staircase, not an artifact of the read-out. The colored/neutral
  confinement split (ratio $\sim 10^3$) is intact.
  \item Each merge step is a genuine valid manifold ($b_1=2$, every triangle in
  $\le 2$ tetrahedra), so the failure to reach the singlet is not a degeneration
  of the cobordism but a property of the two-body color content.
\end{itemize}

\section{Resolution: the trivalent junction \texorpdfstring{$W_{ABC}$}{W\_ABC}}

The representation theory dictates the fix. Since $\varepsilon_{ijk}$ is a
\emph{three}-valent invariant, the proton must be realized by a cobordism with
\emph{three} boundary inputs meeting a single bulk,
\[
  \partial W_{ABC}=\psi_A\sqcup\psi_B\sqcup\psi_C ,
\]
a trivalent junction whose induced orientation on the three color cycles carries
the alternating (sign) structure --- the simplicial $\varepsilon_{ijk}$. This is
the $n=3$ object in which the singlet first occurs, built once rather than as a
sequence of two-body fusions. The exact period machinery used here
(\texttt{readoutHoles}/\texttt{residualForPeriods} to pin the three inputs,
\texttt{cyclePeriods} to read the emergent result, \texttt{emergesResult} to pin
inputs only) is exactly what that junction consumes; only the junction topology
itself is new (tracked as \#396).

\section{Conclusion}

The proton does not emerge from bipartite merges because the $SU(3)$ color
singlet is the three-index alternating invariant $\varepsilon_{ijk}$, which is
absent from $\rep 3\otimes\rep 3=\rep 6\oplus\bar{\rep 3}$ and first appears in
$\rep 3\otimes\rep 3\otimes\rep 3$. The two-quark color Hilbert space simply does
not contain the state we were relaxing toward; a sequence of generic pairwise
fusions cannot manufacture the simultaneous antisymmetry of three color indices.
The qudit form of the statement is exact --- baryon arity $=d$, so qubits make
baryons bipartitely but qutrit color does not --- and our simplicial register
reproduces it quantitatively: the emergent bipartite result stays color-charged
($|\sigmacol|\approx0.67$) at every step. The proton is irreducibly tripartite,
and its carrier is the trivalent junction $W_{ABC}$.

\end{document}
